\exercisesection{挠积}

\begin{enumerate}
\def\labelenumi{\arabic{enumi})}
\item
  我们用 B 表示环 A(ε)（X，第 27 页）；设 C 为 (B, B)-双模的复形，使得 \(C_{n} = B\) 对所有 \(n \in \mathbf{Z}\) , \(d(x) = \varepsilon x\) 对所有 \(x \in C\) 。证明复形 C 的同调为零，但 H \(_{n}\) (C ⊗ \(_{B}\) C) 对所有 \(n \in \mathbf{Z}\) 都同构于 A。由此推出，在引理 1 和命题 4（第 66-67 页）中不能去掉“E 上有界”这一假设。
\item
  设 \(\rho: A \to B\) 是环同态，M 是 A-模。证明以下性质等价：
\end{enumerate}

α) 对每个右 B-模 P，有 \(\operatorname{Tor}_1^{\mathbf{A}}(\rho_*(\mathbf{P}), \mathbf{M}) = 0\) 。

β) 对每个单生成右 B-模 P，有 \(\operatorname{Tor}_1^{\mathbf{A}}(\rho_{*}(\mathbf{P}),\mathbf{M}) = 0\) 。

\(\gamma)\) B-模 \(\mathbf{B} \otimes_{\mathbf{A}} \mathbf{M}\) 是平坦的，且有 \(\operatorname{Tor}_{1}^{\mathbf{A}}(\mathbf{B}, \mathbf{M}) = 0\) .

(为了看出 \(\gamma\) ）蕴含 \(\alpha\) ），考虑右 B-模的正合序列 \(0 \to R \to L \to P \to 0\) ，其中 L 是自由的。）

\begin{enumerate}
\def\labelenumi{\arabic{enumi})}
\setcounter{enumi}{2}
\tightlist
\item
  假设环 A 是交换的。设 I 是一个有限集合，且 \((\mathbf{C}^{(i)}, d^{(i)})_{i \in I}\) 是 A-模复形的一个族。
\end{enumerate}

\begin{enumerate}
\def\labelenumi{\alph{enumi})}
\tightlist
\item
  我们赋予张量积 \(\otimes\) \(\mathbf{C}^{(i)}\) 以其 \(\mathbf{Z}^1\) 型分次以及 A-自同态
\end{enumerate}

\[
d _ {i} = \bigotimes_ {j \in I} d _ {i} ^ {(j)} \quad \text { with } \quad d _ {i} ^ {(j)} = 1 _ {\mathrm{C} ^ {(j)}} \quad \text { for } \quad j \neq i \quad \text { and } \quad d _ {i} ^ {(i)} = d ^ {(i)}.
\]

证明如此便得到一个 I-预复形（X，练习 13，第 174 页）；与之相伴的 I-复形称为复形 \(\mathbf{C}^{(i)}\) 的张量积 I-复形。与该 I-复形相伴的复形称为 \(\mathbf{C}^{(i)}\) 的张量积复形。

\begin{enumerate}
\def\labelenumi{\alph{enumi})}
\setcounter{enumi}{1}
\tightlist
\item
  证明在 I 上选取一个全序，便可定义从 \(\mathbf{C}^{(i)}\) 的张量积复形到 X 第 64 页所定义的复形的一个同构。
\end{enumerate}

\begin{enumerate}
\def\labelenumi{\arabic{enumi})}
\setcounter{enumi}{3}
\tightlist
\item
  设 a 是 A 的一个双边理想。设 M 是一个 A-模，使得 \(\operatorname{Tor}_1^A (\mathbf{A} / \mathfrak{a},\mathbf{M}) = 0\) 并且该 \((\mathbf{A} / \mathfrak{a})\) -模 \(\mathbf{M} / \mathfrak{a}\mathbf{M}\) 是投射的。
\end{enumerate}

\begin{enumerate}
\def\labelenumi{\alph{enumi})}
\item
  假设 M 是有限表示的，且 a 包含在 A 的根中，或者 a 是幂零的。证明 M 是投射的（使用 VIII，§ 8，n° 3，推论 3）。
\item
  假设 M 是有限生成的，且有 \(\bigcap \mathfrak{a}^n = 0\)，并且 (A/a)-模 M/aM 是自由的。证明
\end{enumerate}

A-模 M 是自由的。

\begin{enumerate}
\def\labelenumi{\alph{enumi})}
\setcounter{enumi}{2}
\tightlist
\item
  给出一个局部环 A、其极大理想 m，以及一个有限生成的 A-模 M 的例子，使得 M 不是平坦的，且满足 \(\operatorname{Tor}_{1}^{A}(A/m, M) = 0\) 。
\end{enumerate}

\begin{enumerate}
\def\labelenumi{\arabic{enumi})}
\setcounter{enumi}{4}
\tightlist
\item
  假设存在（含幺）\(k\) -代数的一个同态 \(\pi: \mathbf{A} \to k\) ，它赋予 \(k\) 一个 A-模结构。我们考虑标准分解 \(\mathbf{B}(\mathbf{A}, k)\) （X，第 58 页）；对于 \(n \geqslant 0\) ，我们把 \(\mathbf{B}_n(\mathbf{A}, k)\) 与 \(\mathbf{A} \otimes_k \mathbf{A}^{\otimes^n}\) 等同起来，并用 \(a[a_1, ..., a_n]\) 表示 \(a \otimes a_1 \otimes \cdots \otimes a_n\) ，即 \(\mathbf{A} \otimes_k \mathbf{A}^{\otimes^n}\) 中的元素。
\end{enumerate}

\begin{enumerate}
\def\labelenumi{\alph{enumi})}
\tightlist
\item
  证明复形 \(\mathbf{B}(\mathbf{A}, k) \otimes_k \mathbf{B}(\mathbf{A}, k)\) 定义了 \((\mathbf{A} \otimes_k \mathbf{A})\) -模 \(k\) 的一个分解。证明同态 \((\mathbf{A} \otimes_k \mathbf{A})\) -线性
\end{enumerate}

\[
g: \mathbf {B} (\mathbf {A} \otimes_ {k} \mathbf {A}, k) \rightarrow \mathbf {B} (\mathbf {A}, k) \otimes_ {k} \mathbf {B} (\mathbf {A}, k)
\]

使得

\[
g \left(\left[ x _ {1} \otimes y _ {1}, \dots , x _ {n} \otimes y _ {n} \right]\right) = \sum_ {0 \leqslant p \leqslant n} \pi \left(x _ {p + 1} \dots x _ {n}\right) \left[ x _ {1}, \dots , x _ {p} \right] \otimes y _ {1} \dots y _ {p} \left[ y _ {p + 1}, \dots , y _ {n} \right]
\]

当 \(x_{1},\ldots,x_{n},y_{1},\ldots,y_{n}\) 属于 A 时，是分解的一个态射。

\begin{enumerate}
\def\labelenumi{\alph{enumi})}
\setcounter{enumi}{1}
\tightlist
\item
  假设A是交换的。证明在B(A, k)上存在一个分次A-代数的结构，使得对于 \(x_{1}, \ldots, x_{n} \in A\) :
\end{enumerate}

\[
\left[ x _ {1}, \dots , x _ {p} \right] \left[ x _ {p + 1}, \dots , x _ {n} \right] = \sum_ {\sigma \in M _ {p, n}} \varepsilon (\sigma) \left[ x _ {\sigma (1)}, \dots , x _ {\sigma (n)} \right]
\]

其中 \(M_{p,n}\) 是 \(S_{n}\) 中由这样的置换 \(\sigma\) 所构成的子集： \(\sigma(1) < \cdots < \sigma(p)\) 与 \(\sigma(p + 1) < \cdots < \sigma(n)\) 。

证明B(A, k)于是成为一个微分分次A-代数（X，第183页，练习18）。

\begin{enumerate}
\def\labelenumi{\arabic{enumi})}
\setcounter{enumi}{5}
\tightlist
\item
  设 P 为右 A-模的平坦复形。对每个 A-模 M 及每个整数 i，我们令
\end{enumerate}

\[
\mathrm{T} _ {i} ^ {\mathbf {P}} (\mathrm{M}) = \mathrm{H} _ {i} (\mathrm{P} \otimes_ {\mathbf {A}} \mathrm{M});
\]

如果 \(u:M\to N\) 是一个 A-同态，我们用 \(\mathbf{T}_{i}^{\mathbb{P}}(u)\) 表示 \(k\) -同态 \(\mathbf{H}_{i}(1_{\mathbb{P}}\otimes u)\) 。

\begin{enumerate}
\def\labelenumi{\alph{enumi})}
\tightlist
\item
  设r为整数。证明以下条件等价：
\end{enumerate}

α) 对每个单射 \(u: \mathbf{M}' \to \mathbf{M}\) ，映射 \(\mathbf{T}_{r}^{\mathbb{P}}(u)\) 是单射。

β) 对每个满射 \(v: \mathbf{M} \to \mathbf{M}''\) ，映射 \(\mathbf{T}_{r+1}^{\mathbf{P}}(v)\) 是满射。

\(\gamma)\) 右 A-模 \(\mathbf{P}_{r}/\mathbf{B}_{r}(\mathbf{P})\) 是平坦的。

δ) 存在一个右 A-模的平坦复形 P'，其微分 \(d_{r+1}\) 为零，并且对每个 A-模 M 和每个整数 i，存在 k-同构 \(\varphi_{i}^{\mathrm{M}} : \mathrm{T}_{i}^{\mathrm{P}}(\mathrm{M}) \to \mathrm{T}_{i}^{\mathrm{P}'}(\mathrm{M})\) ，使得

\[
\mathrm{T} _ {i} ^ {\mathrm{P} ^ {\prime}} (u) \circ \varphi_ {i} ^ {\mathrm{M}} = \varphi_ {i} ^ {\mathrm{N}} \circ \mathrm{T} _ {i} ^ {\mathrm{P}} (u)
\]

对每个A-同态 \(u:M\to N\)

（为了证明 \(\gamma\) ）蕴含 \(\delta\) ），定义 \(\mathbf{P}'\) 如下：令 \(\mathbf{P}_{r+1}' = \mathbf{Z}_{r+1}(\mathbf{P})\) ， \(\mathbf{P}_r' = \mathbf{P}_r / \mathbf{B}_r(\mathbf{P})\) 。）

\begin{enumerate}
\def\labelenumi{\alph{enumi})}
\setcounter{enumi}{1}
\tightlist
\item
  我们还假设A是诺特环，且A-模 \(P_{i}\) 是有限生成的。证明a)中的条件等价于以下条件：
\end{enumerate}

e) 存在一个A-模Q，并且对每个A-模M，有一个k-同构 \(\psi^{M}: T_{r}^{P}(M) \to \text{Hom}_{A}(Q, M)\) 使得 \(\psi^{N} \circ T_{r}^{P}(u) = \text{Hom}(1, u) \circ \psi^{M}\) 对每个A-同态 \(u: M \to N\) .

(为证明δ)蕴含ε)，可化简为以下情形： \(d_{r+1} = 0\) 并取Q为 \(Z_{r}(\mathbf{P})\) 的对偶。) c) 证明在A是诺特环的假设下，b)的结果仍然成立， \(H_{i}(\mathbf{P})\) 对每个 i 都是有限生成的，并且 \(H_{i}(\mathbf{P}) = 0\) 当 \(i < i_{0}\) （利用命题 10，第 56 页）。

\begin{enumerate}
\def\labelenumi{\arabic{enumi})}
\setcounter{enumi}{6}
\tightlist
\item
  设 C 是右 A-模的复形，M 是左 A-模。C 与 M 的挠积是分次 k-模 \(\mathcal{T}or^{\mathbf{A}}(\mathbf{C},\mathbf{M})=\mathrm{H}(\mathbf{C}\otimes_{\mathbf{A}}\mathbf{L}(\mathbf{M}))\) 如果 \(f:C\to C'\) 是一个复形的态射且 \(g:M\to M'\) 一个A-模的同态，我们记 \(\mathcal{T}or^{\mathbf{A}}(f,g)=\mathrm{H}(f\otimes\mathbf{L}(g))\) 。 a) 如果 \(f:C\to C'\) 是一个拟同构， \(\mathcal{T}or^{\mathbf{A}}(f,1)\) 是一个同构；如果C是平坦的且在右边有界，则分次k-模 \(\mathcal{T}or^{\mathbf{A}}(\mathbf{C},\mathbf{M})\) 同构于 \(\mathrm{H}(\mathbf{C}\otimes_{\mathbf{A}}\mathbf{M})\) 。
\end{enumerate}

\begin{enumerate}
\def\labelenumi{\alph{enumi})}
\setcounter{enumi}{1}
\tightlist
\item
  若 \(0 \to C' \to C \to C'' \to 0\) 是右A-模复形的一个正合序列，且若M是一个A-模，证明存在一个正合序列
\end{enumerate}

\[
\dots \rightarrow \mathcal {T} o r _ {n} ^ {\mathrm{A}} (\mathrm{C} ^ {\prime}, \mathrm{M}) \rightarrow \mathcal {T} o r _ {n} ^ {\mathrm{A}} (\mathrm{C}, \mathrm{M}) \rightarrow \mathcal {T} o r _ {n} ^ {\mathrm{A}} (\mathrm{C} ^ {\prime \prime}, \mathrm{M}) \rightarrow \mathcal {T} o r _ {n - 1} ^ {\mathrm{A}} (\mathrm{C} ^ {\prime}, \mathrm{M}) \rightarrow \dots .
\]

若 \(0 \to M' \xrightarrow{i} M \xrightarrow{p} M'' \to 0\) 是A-模的一个正合序列，证明存在一个正合序列：

\[
\dots \rightarrow \mathcal {T} o r _ {n} ^ {\mathrm{A}} (\mathrm{C}, \mathrm{M} ^ {\prime}) \rightarrow \mathcal {T} o r _ {n} ^ {\mathrm{A}} (\mathrm{C}, \mathrm{M}) \rightarrow \mathcal {T} o r _ {n} ^ {\mathrm{A}} (\mathrm{C}, \mathrm{M} ^ {\prime \prime}) \rightarrow \mathcal {T} o r _ {n - 1} ^ {\mathrm{A}} (\mathrm{C}, \mathrm{M} ^ {\prime}) \rightarrow \dots .
\]

（注意存在一个从L(M'')到L(i)的锥的同伦映射）。

\begin{enumerate}
\def\labelenumi{\alph{enumi})}
\setcounter{enumi}{2}
\tightlist
\item
  证明存在一个谱序列 'E（X，第175–177页，习题14–17）收敛到 \(\mathcal{T}or^{\mathrm{A}}(\mathrm{C},\mathrm{M})\) ，使得 \(E_{pq}^{2} = \mathrm{Tor}_{p}^{\mathrm{A}}(\mathrm{H}_{q}(\mathrm{C}),\mathrm{M})\) 此外，若复形 C 右有界，证明存在一个谱序列 E 收敛到 \(\mathcal{T}or^{\mathrm{A}}(\mathrm{C},\mathrm{M})\) ，使得 \(E_{pq}^{2} = H_{p}(\mathrm{Tor}_{q}^{\mathrm{A}}(\mathrm{C},\mathrm{M}))\) （考虑双复形 \(\mathrm{C}\otimes_{\mathrm{A}}\mathrm{L}(\mathrm{M})\) ）。
\end{enumerate}

\begin{enumerate}
\def\labelenumi{\arabic{enumi})}
\setcounter{enumi}{7}
\tightlist
\item
  设A、B为两个k-代数（结合的且含幺），M为右A-模，P为左B-模，N为(A, B)-双模。
\end{enumerate}

\begin{enumerate}
\def\labelenumi{\alph{enumi})}
\tightlist
\item
  证明存在两个 \(k\) -模的谱序列 \({}^{1}\mathrm{E}\) 和 \({}^{2}\mathrm{E}\) ，收敛到相同的 \(k\) -分次模，使得
\end{enumerate}

\[
{ } ^ { 1 } \mathrm{E} _ { p q } ^ { 2 } = \operatorname{Tor} _ { p } ^ { \mathrm{A} } ( \mathrm{M} , \operatorname{Tor} _ { q } ^ { \mathrm{B} } ( \mathrm{N} , \mathrm{P} ) ) ; \quad { } ^ { 2 } \mathrm{E} _ { p q } ^ { 2 } = \operatorname{Tor} _ { p } ^ { \mathrm{B} } ( \operatorname{Tor} _ { q } ^ { \mathrm{A} } ( \mathrm{M} , \mathrm{N} ) , \mathrm{P} ) .
\]

（考虑双复形 \(L(M) \otimes_{A} (N \otimes_{B} L(P))\) .)

\begin{enumerate}
\def\labelenumi{\alph{enumi})}
\setcounter{enumi}{1}
\item
  如果 N 是平坦的 B-模，证明谱序列 \({}^{2}E\) 收敛到 Tor \(^{A}\) (M, N ⊗ \(_{B}\) P)。如果 N 是平坦的 A-模，序列 \({}^{1}E\) 收敛到 Tor \(^{B}\) (M ⊗ \(_{A}\) N, P)。
\item
  设 \(\rho: \mathbf{A} \to \mathbf{B}\) 为 \(k\) -代数的一个同态。由 \(b\) ）推出一个谱序列 E，使得
\end{enumerate}

\[
\mathrm{E} _ {p q} ^ {2} = \operatorname{Tor} _ {p} ^ {\mathbf {B}} \left(\operatorname{Tor} _ {q} ^ {\mathbf {A}} (\mathbf {M}, \mathbf {B}), \mathbf {P}\right),
\]

收敛到 Tor \(^{A}\) (M, P)。

\begin{enumerate}
\def\labelenumi{\arabic{enumi})}
\setcounter{enumi}{8}
\tightlist
\item
  设 P 是右 A-模的复形，Q 是左 A-模的复形。假设 P 和 Q 都下有界，或者存在一个整数 d，使得每个左或右 A-模都有长度为 \(\leqslant d\) 的投射分解。
\end{enumerate}

\begin{enumerate}
\def\labelenumi{\alph{enumi})}
\tightlist
\item
  证明存在两个谱序列（X，第 175 页，习题 14）'E 和 ``E，收敛到同一个分次 k-模，使得
\end{enumerate}

\[
^ {\prime} \mathrm{E} _ {p q} ^ {2} = \mathrm{H} _ {p} (\operatorname{Tor} _ {q} ^ {\mathbf {A}} (\mathrm{P}, \mathrm{Q})); \quad^ {\prime \prime} \mathrm{E} _ {p q} ^ {2} = \bigoplus_ {q ^ {\prime} + q ^ {\prime \prime} = q} \operatorname{Tor} _ {p} ^ {\mathbf {A}} \left(\mathrm{H} _ {q ^ {\prime}} (\mathrm{P}), \mathrm{H} _ {q ^ {\prime \prime}} (\mathrm{Q})\right)
\]

\((\mathrm{where}\ \mathrm{Tor}_{q}^{\mathrm{A}}(\mathrm{P},\mathrm{Q})\ \mathrm{denotes}\ \mathrm{the}\ \mathrm{complex}\ \mathrm{associated}\ \mathrm{with}\ \mathrm{the}\ \mathrm{bicomplexe}\ \oplus\ \mathrm{Tor}_{q}^{\mathrm{A}}(\mathrm{P}_{r},\mathrm{Q}_{s}))\)

（考虑 P 和 Q 的投射分解 \(\mathcal{P}\) 和 \(\mathcal{Q}\) （X，第 183 页，习题 17），以及满足 \(C_{p,q} = \bigoplus (\mathcal{P}_{p'q'} \otimes \mathcal{Q}_{p''q''})\) 的双复形 C。）

\begin{enumerate}
\def\labelenumi{\alph{enumi})}
\setcounter{enumi}{1}
\tightlist
\item
  如果 P 是平坦的，证明谱序列 ``E 收敛到 H(P ⊗\_A Q)。如果 H(P) 是平坦的，序列 'E 收敛到 H(P) ⊗\_A H(Q)。
\end{enumerate}

